\documentclass[11pt]{article}
\usepackage[margin=1in]{geometry}
\usepackage{amsmath,amssymb,amsthm}
\usepackage[T1]{fontenc}
\usepackage{lmodern}
\usepackage{microtype}
\usepackage[colorlinks=true,linkcolor=blue,urlcolor=blue]{hyperref}
\hypersetup{pdftitle={A negative solution to MAIS-O87},pdfauthor={Robert Sneiderman}}
\newtheorem{theorem}{Theorem}
\newtheorem{lemma}{Lemma}
\newcommand{\eps}{\varepsilon}
\newcommand{\kap}{\kappa}
\newcommand{\R}{\mathbb R}
\newcommand{\grad}{\nabla}
\title{A negative solution to MAIS-O87}
\author{Robert Sneiderman\\Independent\\\texttt{robbysneiderman@gmail.com}}
\date{September 2026}
\begin{document}
\maketitle
\begin{abstract}
In the coin-line model of MAIS-O87, the training coin sits at the right
endpoint in most episodes, and the question is whether the learned preference
for moving right can survive at a test state whose coin lies to the left.
It cannot. For every fixed positive training diversity and every finite
initialization, linear policy-gradient flow eventually assigns probability
less than one half to moving right at the test state. The misgeneralization
probability therefore tends to zero, which answers O87 in the negative. The
correction time may depend on the diversity and on the initial parameters;
no rate is claimed.
\end{abstract}

\section{The question and the result}
The agent moves on the sites $0,\ldots,L$ and receives reward for reaching
a coin. During training, the coin is usually at $L$, but with probability
$\eps$ it is placed uniformly on $1,\ldots,L$. A linear logistic policy
chooses its direction from the agent's position and the coin's position.
At the test state, the agent starts at $\lceil L/2\rceil$ and the coin is
at zero, so moving left is the correct action.

MAIS-O87~\cite{mais-o87} asks whether, after arbitrarily long training, a positive
fraction of Gaussian initializations can still favor moving right at this
test state, uniformly over small positive $\eps$. The following theorem
rules this out. The coordinates $(a,b,d)$ and the gradient flow are specified
in the next section; $q^{\mathrm{RL}}_\eps(t;\sigma,L)$ denotes the probability
that the test-state logit is positive at time $t$ when
$w(0)\sim N(0,\sigma^2 I_3)$.

\begin{theorem}\label{thm:main}
Fix an integer $L\ge4$ and a real $\eps\in(0,1]$. In the MAIS-A8
coin-line model, every finite initialization of linear policy-gradient flow
satisfies
\[
 a(t)+\lceil L/2\rceil b(t)<0\qquad\text{for all sufficiently large }t.
\]
Consequently, for every $\sigma>0$,
\[
 \lim_{t\to\infty}q^{\mathrm{RL}}_\eps(t;\sigma,L)=0.
\]
Thus the positive uniform lower bound asked for in O87 does not exist.
\end{theorem}
The conclusion is pointwise in $\eps$. Each fixed positive value gives
eventual correction, with no bound on the correction time. That is enough
for O87, which takes the long-time limit with $\eps$ held fixed.

\paragraph{Idea of the proof.}
The parameter difference $A=d-b$ increases without bound, while every
individual logit derivative tends to zero. At a state next to a coin, this
forces the logit to tend to either $+\infty$ or $-\infty$. The main step is
to rule out the possibility that all these adjacent states eventually favor
moving right.

The reason is an asymmetry between the two sides of a coin. When the
logits below the coin are positive, an episode that starts below and fails
must make at least $\lceil(L+1)/2\rceil$ left moves, each unlikely, so the
gradient from that side decays like a high power of a small probability.
Just above the coin, the adjacent state can refuse the coin by never stepping
left, and the derivative there carries a single small factor, the probability
of stepping left; every other factor is bounded below independently of the
rest of the policy. The corrective gradient from above therefore decays more
slowly than the gradient from below, and any positive training weight on
coins left of $L$ is enough to break the all-right pattern. Monotone bounds
on the parameters then show that a nonnegative test-state logit at
arbitrarily late times would force exactly that pattern. The proof does not
need the training return to converge to one.

\section{The model and its gradient}
Write $w=(a,b,d)$ and $H=2L$. At position $p$ with coin at $c$, the
logit and the probability of moving right are
\[
 z_{pc}=a+bp+dc,\qquad r_{pc}=\frac1{1+e^{-z_{pc}}}.
\]
The position is $p\in\{0,\ldots,L\}$; the training coin is
$c\in\{1,\ldots,L\}$, with probability
\[
 \mu_c=\frac{\eps}{L}+(1-\eps)\mathbf1_{\{c=L\}}>0.
\]
Conditional on $c$, the initial position is uniform on $\{0,\ldots,L\}$.
An episode stops successfully on hitting $c$, including at time zero, and
otherwise stops unsuccessfully after $H$ steps. Left and right steps are
clipped to $[0,L]$. The time-zero success and $H$-action conventions are
those of MAIS-A8, Definition 3.1~\cite{mais-a8}.
The objective $J$ is success probability, and
$\dot w=\grad J$ is Euclidean gradient ascent in these three coordinates.

Let $J_c$ be conditional success probability. Temporarily regard all
nonterminal logits for fixed $c$ as independent variables, and define
\[
 y_{pc}=\operatorname{sign}(c-p),\qquad
 \kap_{pc}=\left|\frac{\partial J_c}{\partial z_{pc}}\right|.
\]
Increasing a logit helps below the coin and hurts above it. More precisely,
\begin{align}
 \frac{\partial J_c}{\partial z_{pc}}&=y_{pc}\kap_{pc},
       &\kap_{pc}&>0,\label{eq:sign}\\
 \grad J&=\sum_{c=1}^L\sum_{p\ne c}
       \mu_c\kap_{pc}y_{pc}(1,p,c).\label{eq:gradient}
\end{align}
To prove \eqref{eq:sign}, let $V_h(p,c)$ denote probability of hitting $c$
within $h$ steps, and let $\rho_t(p,c)$ be the probability, conditional on $c$,
of occupying $p$ at time $t$ without having stopped. For fixed $c$, the Bellman recursion is
\[
 V_h(s,c)=(1-r_{sc})V_{h-1}(s^-,c)+r_{sc}V_{h-1}(s^+,c)
 \quad(s\ne c),
\]
with $V_0(s,c)=\mathbf1_{\{s=c\}}$ and $V_h(c,c)=1$.
Fix a nonterminal $p$ and put $D_h(s)=\partial V_h(s,c)/\partial z_{pc}$.
Differentiating one step yields
\begin{align*}
 D_h(s)={}&(1-r_{sc})D_{h-1}(s^-)+r_{sc}D_{h-1}(s^+)\\
 &+\mathbf1_{\{s=p\}}r_{pc}(1-r_{pc})
   [V_{h-1}(p^+,c)-V_{h-1}(p^-,c)].
\end{align*}
Here $D_0=0$ and $D_h(c)=0$. Iterating this identity and averaging the initial
state gives the occupancy weights $\rho_t$ and hence
\begin{equation}\label{eq:occupancy}
 \frac{\partial J_c}{\partial z_{pc}}
 =r_{pc}(1-r_{pc})\sum_{t=0}^{H-1}\rho_t(p,c)
 \bigl[V_{H-t-1}(p^+,c)-V_{H-t-1}(p^-,c)\bigr],
\end{equation}
where $p^+=\min(p+1,L)$ and $p^-=\max(p-1,0)$.

Below the coin, a path from $p^-$ must hit $p^+$ before reaching $c$.
The strong Markov property and monotonicity in the remaining horizon give
$V_h(p^-,c)\le V_h(p^+,c)$. Reflection gives the opposite inequality above
the coin. At $h=H-1$ the inequality is strict: the closer neighbor has positive
success probability, whereas from the farther neighbor there is positive
probability of avoiding the closer one for all $h$ steps. This uses self-loops
at a wall or back-and-forth steps away from the closer neighbor. Since
$\rho_0(p,c)=1/(L+1)$, equation \eqref{eq:occupancy} has the strict claimed
sign. Equation \eqref{eq:gradient} is the chain rule.

\section{Three monotone parameter combinations}
The following combinations control how the parameters can diverge:
\[
 A=d-b,\qquad R=A-a,\qquad S=A+a.
\]
The flow is globally defined. Indeed, $J$ is a finite polynomial in finitely
many logistic probabilities of linear forms. Its gradient and Hessian are
globally bounded, so the gradient field is globally Lipschitz and bounded.
Equation \eqref{eq:gradient} gives
\begin{equation}\label{eq:escape}
 A'=\sum_{c,p\ne c}\mu_c\kap_{pc}|c-p|>0,
 \qquad \|w'\|\le\sqrt{1+2L^2}\,A'.
\end{equation}
If $A$ had a finite upper bound, \eqref{eq:escape} would imply finite total
arclength and convergence of $w$ to a finite point. At that point $A'$ is
strictly positive, a contradiction. Therefore
\begin{equation}\label{eq:Ainfty} A(t)\longrightarrow+\infty.\end{equation}
Also $J'=\|\grad J\|^2$ and $J\le1$. Thus the squared gradient has finite
integral. It is uniformly continuous, since the gradient and Hessian are
bounded. A nonnegative, uniformly continuous, integrable function tends to
zero; hence $\grad J\to0$. Equation \eqref{eq:escape} implies $A'\to0$ and,
because every $\mu_c$ is positive,
\begin{equation}\label{eq:kzero} \kap_{pc}(t)\longrightarrow0\quad(p\ne c).\end{equation}
Finally,
\begin{align}
 R'&=\sum_{c,p\ne c}\mu_c\kap_{pc}(|c-p|-y_{pc})\ge0,\label{eq:R}\\
 S'&=\sum_{c,p\ne c}\mu_c\kap_{pc}(|c-p|+y_{pc})\ge0.\label{eq:S}
\end{align}
Only the lower bounds $R\ge R(0)$ and $S\ge S(0)$, not divergence of $R$ or
$S$, are needed below.

\section{The logits next to a coin}
The two states next to each training coin carry the argument. Write
\[
 \ell_c=z_{c-1,c}\quad(1\le c\le L),\qquad
 q_c=z_{c+1,c}\quad(1\le c<L).
\]
\begin{lemma}\label{lem:signs}
Each adjacent logit tends either to $+\infty$ or to $-\infty$.
\end{lemma}
\begin{proof}
At a left-adjacent state, retain $t=0$ in \eqref{eq:occupancy}. Right
immediately succeeds. After choosing left, success requires a later right
action at $c-1$. Choosing left at every subsequent visit to $c-1$ for $H-1$
steps prevents success, with probability at least $(1-r)^{H-1}$. Therefore
\begin{equation}\label{eq:leftbound}
 \kap_{c-1,c}\ge\frac{r(1-r)^H}{L+1}.
\end{equation}
Similarly, at a right-adjacent state,
\begin{equation}\label{eq:rightbound}
 \kap_{c+1,c}\ge\frac{r^H(1-r)}{L+1}.
\end{equation}
These estimates are independent of all other policy probabilities. Precisely,
modify the policy at the adjacent state to choose deterministically away from
the coin. Under the modified policy, hitting the coin from the away neighbor
is impossible. The original probability of these paths is the expectation of
$r^N$, or $(1-r)^N$, with $N\le H-1$. This also includes clipped wall actions.

On a compact interval of logits the right sides of
\eqref{eq:leftbound}--\eqref{eq:rightbound} have positive lower bounds.
Thus \eqref{eq:kzero} makes the absolute value of each adjacent logit tend to
infinity. Continuity then gives a fixed eventual sign.
\end{proof}

\section{The adjacent logits cannot all tend to positive infinity}
\begin{lemma}\label{lem:exclude}
No trajectory has $\ell_c\to+\infty$ for every $c$ and $q_c\to+\infty$
for every $c<L$.
\end{lemma}
The lemma places no condition on logits farther from the coin, so the
estimates below must hold whatever the policy does at those states.

\subsection{Gradient estimates}
Write $\grad J_c=P_c-N_c$, where $P_c$ is the contribution from $p<c$ and
$N_c$ is the magnitude of the contribution from $p>c$. Both are componentwise
nonnegative by \eqref{eq:sign}; $N_L=0$.
Let $F_c(p_0)=1-V_H(p_0,c)$. A path starting below the coin stops before
visiting any state above it, so
\[
 P_c=\frac1{L+1}\sum_{p_0<c}\grad V_H(p_0,c)
    =-\frac1{L+1}\sum_{p_0<c}\grad F_c(p_0).
\]
For the finite set $\mathcal F_c(p_0)$ of failed length-$H$ action paths,
\[
 \grad F_c(p_0)=\sum_{\omega\in\mathcal F_c(p_0)}
 \Pr_w(\omega\mid p_0,c)\grad\log\Pr_w(\omega\mid p_0,c).
\]
Put $n=\lceil(L+1)/2\rceil\ge3$. Constants $C,K>0$ depending only on $L$
satisfy the following estimates for every coordinate $j$.
If $m_c=\min_{p<c}z_{pc}\ge0$, then
\begin{equation}\label{eq:Pbound} (P_c)_j\le C e^{-nm_c}.\end{equation}
For $c=1$ and $h=z_{01}\ge0$,
\begin{equation}\label{eq:Pone}
 (P_1)_a=(P_1)_d\le Ce^{-Hh},\qquad (P_1)_b=0.
\end{equation}
If $q_c\ge0$ and $c<L$, then, without restrictions on other logits,
\begin{equation}\label{eq:Nbound} (N_c)_j\ge K e^{-q_c}.\end{equation}
One may take $C=HL2^H$ and $K=2^{-(H+1)}/(L+1)$.

For \eqref{eq:Pbound}, a failed path starting below $c$ stays below $c$.
If it makes $N_-$ left moves, its final position is at least $p_0+H-2N_-$;
clipping at zero only increases this lower bound. Consequently
$N_-\ge(H+p_0-c+1)/2$, and hence $N_-\ge n$. Each left probability is at
most $e^{-m_c}$, so $\Pr_w(\omega\mid p_0,c)\le e^{-nm_c}$ when $m_c\ge0$.
Each of the $H$ factors of $\log\Pr_w(\omega\mid p_0,c)$ has partial
derivatives bounded by $L$ in absolute value, so the score is at most $HL$
in every coordinate. There are at most $2^H$ action paths and at most $L$
starting positions, each of weight $1/(L+1)$, and \eqref{eq:Pbound} follows
with $C=HL2^H$.
For $c=1$, failure from
below requires $H$ left self-loops at $p=0$, proving \eqref{eq:Pone}.
For \eqref{eq:Nbound}, use \eqref{eq:rightbound}, $r\ge1/2$, and
$r(1-r)\ge e^{-q_c}/4$. All components of $(1,c+1,c)$ are at least one.

\subsection{Proof of the exclusion lemma}
\begin{proof}[Proof of Lemma \ref{lem:exclude}]
Suppose all adjacent logits tend to $+\infty$. At sufficiently late times
$A>0$ and $h=\ell_1>0$. Whenever $b\ge0$, we have $d=A+b>0$ and
\[
 m_c=h+(c-1)d,\qquad q_1=h+2b.
\]
Only $c\ge2$ contribute positively to $b'$. With $\mu_1=\eps/L$,
\begin{equation}\label{eq:bdrift}
 b'\le Ce^{-n(h+d)}-\mu_1Ke^{-(h+2b)}.
\end{equation}
The ratio of the positive term to the negative term is at most
\[
 \frac{C}{\mu_1K}\exp\bigl(-(n-1)h-(n-2)b-nA\bigr),
\]
which tends uniformly to zero as $A\to\infty$ for $h,b\ge0$. Thus, at
sufficiently late times, $b'<0$ whenever $b\ge0$. A first late crossing from $b<0$ to $b\ge0$ would occur at $b=0$ with
nonnegative left difference quotient, contradicting $b'<0$ there. Hence
no such crossing is possible.

Two cases remain. If $b$ stays nonnegative, it is bounded and decreasing,
while $d\to\infty$ and $h\to\infty$ by hypothesis. The positive $a$ and $d$ components are bounded by
\[
 Ce^{-Hh}+Ce^{-n(h+d)},
\]
while their negative contribution from $c=1$ is at least
$\mu_1Ke^{-(h+2b)}$. Since $b$ is bounded, both positive-to-negative ratios
tend to zero. Thus $a',d'<0$ eventually, contradicting $h=a+d\to\infty$.

Otherwise $b<0$ eventually. For $c<L$,
\[
 m_c=\ell_c=q_c-2b\ge q_c.
\]
As $q_c\to\infty$, equations \eqref{eq:Pbound} and \eqref{eq:Nbound} give,
in every coordinate,
\begin{equation}\label{eq:negative}
 (\grad J_c)_j\le-\tfrac K2 e^{-q_c}
\end{equation}
eventually. For $c=L$, use the exact identity
\[
 m_L=\ell_L=q_{L-1}+A.
\]
Its positive gradient is at most $Ce^{-n(q_{L-1}+A)}$, which is eventually
smaller than half the weighted negative contribution of coin $L-1$ in
\eqref{eq:negative}. Every component of $\grad J$ is therefore strictly
negative eventually. This gives $q_1'=a'+2b'+d'<0$, contradicting
$q_1\to\infty$. Both cases are impossible.
\end{proof}

\section{From the adjacent states to the test state}
The last step connects the adjacent states to the test state. Let
$k=\lceil L/2\rceil\ge2$ and let $P=a+kb$ be the test-state logit. Along
any sequence of times tending to infinity with $P\ge0$, every adjacent logit
tends to $+\infty$; the two signs of $b$ are treated separately. Write
$R_0=R(0)$ and $S_0=S(0)$.

If $b\ge0$, substituting $a=P-kb$ into $S=A+a\ge S_0$ gives
\[
 b\le\frac{A+P-S_0}{k}.
\]
Consequently
\begin{equation}\label{eq:casepositive}
 h=a+d=P+A-(k-1)b
 \ge\frac{A+P}{k}+\frac{k-1}{k}S_0
 \ge\frac Ak+\frac{k-1}{k}S_0.
\end{equation}
For $A>0$, $d=A+b>0$, and every training logit satisfies
$z_{pc}=h+bp+(c-1)d\ge h$. Equation \eqref{eq:Ainfty} proves the claim
in this case.

If $b<0$, write $B=-b>0$. Then $a=P+kB$, and $R=A-a\ge R_0$ gives
\[
 kB\le A-P-R_0.
\]
Writing $D=b+d=A-2B$, we have
\[
 D=R+P+(k-2)B\ge R_0,\qquad q_1=P+A+(k-3)B.
\]
For $k\ge3$, $q_1\ge A$. For $k=2$,
\[
 q_1\ge\frac A2+\frac{3P}{2}+\frac{R_0}{2}
       \ge\frac A2+\frac{R_0}{2}.
\]
At sufficiently late times $A>0$. In either case $q_1\ge A/2-|R_0|/2$. Since $q_c=q_1+(c-1)D$,
\begin{equation}\label{eq:casenegative}
 q_c\ge\frac A2-L|R_0|\qquad(1\le c<L).
\end{equation}
The other adjacent logits satisfy $\ell_c=q_c+2B$ for $c<L$ and
$\ell_L=q_{L-1}+A$. They too diverge positively by \eqref{eq:Ainfty}.
This completes the claim for both signs of $b$.

\section{Proof of the theorem}
\begin{proof}[Proof of Theorem \ref{thm:main}]
If $P(t)\ge0$ at arbitrarily late times, the previous section supplies an
unbounded sequence of times along which all adjacent logits are positive.
Lemma \ref{lem:signs} says that each has a fixed eventual sign and diverges in
magnitude. All must therefore tend to $+\infty$, contradicting
Lemma \ref{lem:exclude}.

Thus $P(t)<0$ for all sufficiently large $t$, from every finite initial
vector. The indicator $\mathbf1_{\{P(t)>0\}}$ tends pointwise to zero and
is bounded by one. Bounded convergence under $N(0,\sigma^2I_3)$ gives
$q^{\mathrm{RL}}_\eps(t;\sigma,L)\to0$. The same conclusion holds for any
probability law on finite initial vectors.
\end{proof}
For each fixed positive $\eps$, the limsup in O87 is therefore zero.
In particular, it cannot have a positive lower bound uniform over small
positive $\eps$.

\section*{Acknowledgments}
Computations and draft review used Astra and Fable 5.1.

\begin{thebibliography}{9}
\bibitem{mais-a8}
MAIS project (\texttt{github.com/lionellevine/MAIS}),
\emph{A predictive theory of out-of-distribution generalization:
which of two perfect policies does gradient descent select?}
Research agenda MAIS-A8, Definition 3.1 and Question 6.2,
\href{https://github.com/lionellevine/MAIS/blob/7b23a8a30c0eb0b1d3b55c592ca9554256fc4805/agendas/A8/MAIS-A8.tex}{pinned source, commit 7b23a8a}.

\bibitem{mais-o87}
MAIS project (\texttt{github.com/lionellevine/MAIS}),
\emph{Does policy gradient stay misgeneralized despite training diversity?}
Open problem MAIS-O87,
\href{https://github.com/lionellevine/MAIS/blob/7b23a8a30c0eb0b1d3b55c592ca9554256fc4805/open-problems/MAIS-O87.md}{problem statement, commit 7b23a8a}.
\end{thebibliography}
\end{document}
